% GB28181协议
% GB28181.tex

\documentclass[12pt,UTF8]{ctexbook}

% 设置纸张信息。
\input{../../../common/page}

% 设置字体，并解决显示难检字问题。
\xeCJKsetup{AutoFallBack=true}
% 注意：ctexbook类已默认设置SimSun为CJKrmdefault
% 以下设置用于确保字体回退和扩展字体可用
% 仅设置扩展字体以避免与默认设置冲突
\setCJKfamilyfont{hei}{SimHei}
\setCJKfamilyfont{kai}{KaiTi}

% 目录 chapter 级别加点（.）。
\usepackage{titletoc}
\titlecontents{chapter}[0pt]{\vspace{3mm}\bf\addvspace{2pt}\filright}{\contentspush{\thecontentslabel\hspace{0.8em}}}{}{\titlerule*[8pt]{.}\contentspage}

% 设置 part 和 chapter 标题格式。
\ctexset{
chapter/name={第,章},
chapter/number={\arabic{chapter}}
}

% 图片相关设置。
\usepackage{graphicx}
\graphicspath{{Images/}}

% 设置署名格式。
\newenvironment{shuming}{\hfill\zihao{4}}

% 注脚每页重新编号，避免编号过大。
\usepackage[perpage]{footmisc}

\title{\heiti\zihao{0} GB28181协议详解}
\author{佚名}
\date{}

\begin{document}

\maketitle
\tableofcontents

\frontmatter

\mainmatter

\chapter{GB28181协议概述}

GB28181是中国国家标准《安全防范视频监控联网系统信息传输、交换、控制技术要求》，是由公安部科技信息化局提出，由全国安全防范报警系统标准化技术委员会(SAC/TC100)归口，公安部第一研究所等单位起草的国家标准。

\section{基本概念}

\begin{itemize}
  \item \textbf{定义}：GB28181是中国公安部主导制定的视频监控联网系统国家标准，规定了城市监控报警联网系统中信息传输、交换、控制的技术要求
  \item \textbf{发布时间}：2011年发布，2016年和2022年分别进行了修订
  \item \textbf{适用范围}：适用于安全防范视频监控联网系统的设计、施工、检验和验收
  \item \textbf{核心目标}：实现不同厂商视频监控设备的互联互通
\end{itemize}

\section{核心功能}

\begin{enumerate}
  \item \textbf{设备注册与管理}：设备的注册、注销、状态查询等
  \item \textbf{实时视频点播}：远程调看监控视频
  \item \textbf{设备控制}：PTZ控制、录像控制等
  \item \textbf{报警事件处理}：报警的上报、分发和处理
  \item \textbf{录像检索与回放}：历史录像的查询和回放
  \item \textbf{设备信息查询}：查询设备的能力和状态
\end{enumerate}

\chapter{技术特点}

\begin{itemize}
  \item \textbf{协议基础}：基于SIP（Session Initiation Protocol）协议扩展
  \item \textbf{传输层}：基于TCP/UDP协议
  \item \textbf{媒体传输}：支持RTP/RTCP协议传输音视频
  \item \textbf{编码支持}：支持H.264、H.265等编码格式
  \item \textbf{安全性}：支持设备认证和媒体流加密
\end{itemize}

\chapter{工作原理}

\begin{enumerate}
  \item \textbf{设备注册}：监控设备向SIP服务器注册
  \item \textbf{设备发现}：客户端通过SIP服务器发现网络中的设备
  \item \textbf{会话建立}：建立视频流传输会话
  \item \textbf{媒体传输}：通过RTP协议传输音视频数据
  \item \textbf{会话控制}：控制视频流的开始、暂停、结束等
  \item \textbf{设备管理}：对设备进行配置和管理
\end{enumerate}

\chapter{与其他标准的对比}

\section{GB28181 vs ONVIF}

\begin{itemize}
  \item \textbf{适用范围}：GB28181主要在中国国内使用，ONVIF是国际标准
  \item \textbf{协议基础}：GB28181基于SIP，ONVIF基于SOAP/HTTP
  \item \textbf{功能特点}：GB28181更注重报警联动和平台级联，ONVIF更注重设备级互操作
  \item \textbf{政府支持}：GB28181是中国国家标准，得到政府强制推行
\end{itemize}

\section{GB28181 vs RTSP}

\begin{itemize}
  \item \textbf{作用范围}：GB28181是完整的监控系统标准，RTSP只是媒体流控制协议
  \item \textbf{协议层次}：GB28181包含设备管理、报警等多种功能，RTSP仅负责媒体流控制
  \item \textbf{应用场景}：GB28181适用于大型监控联网系统，RTSP适用于点对点媒体流控制
\end{itemize}

\chapter{应用场景}

\begin{itemize}
  \item \textbf{平安城市}：城市级视频监控联网系统
  \item \textbf{智能交通}：道路监控、交通管理系统
  \item \textbf{安防监控}：企业、园区、小区的视频监控系统
  \item \textbf{应急指挥}：突发事件的视频调度和指挥
  \item \textbf{雪亮工程}：农村地区视频监控全覆盖工程
\end{itemize}

\chapter{优缺点}

\section{优点}

\begin{itemize}
  \item 符合中国国情，适应国内安防系统需求
  \item 支持大规模联网，适合城市级监控系统
  \item 统一标准，减少系统集成成本
  \item 政府支持，推广力度大
  \item 功能完善，涵盖监控系统各方面需求
\end{itemize}

\section{缺点}

\begin{itemize}
  \item 国际化程度低，主要在中国国内使用
  \item 协议相对复杂，实现难度较大
  \item 与国际标准存在一定差异，需要额外适配
  \item 版本更新较慢，可能跟不上技术发展
\end{itemize}

\chapter{实际应用}

在国内安防项目中，GB28181已成为强制标准。例如，平安城市项目要求所有接入的监控设备必须支持GB28181协议，确保不同厂商的设备能够无缝集成到统一的监控平台中。公安部门通过GB28181协议实现对辖区内所有监控设备的统一管理和调看，提高了应急处置能力。

随着技术的发展，GB28181标准也在不断更新，2022年发布的新版本增加了对AI分析、智能识别等新技术的支持，进一步提升了标准的适用性和先进性。

\backmatter

\end{document}